\documentclass[12pt]{article}
\usepackage[utf8]{inputenc}
\usepackage[T1]{fontenc}
\usepackage{hyperref}
\title{Suena Language Resources for K-12 Curriculum Development}
\author{Compiled by Claude Sonnet 5.5 (claude-sonnet-5-5) for Andr\'as Kornai}
\date{}
\begin{document}
\maketitle

\section{Introduction}
Suena (ISO 639-3: sue; Glottocode suen1241) is a Papuan language of Morobe Province, Papua New Guinea, in the Binanderean branch of Trans--New Guinea. Speaker estimates vary: about 3,600 (2000) in one source, 2,275 in the Global Recordings Network listing. UNESCO classes it as critically endangered. Other names are Yarawe, Yarawi and Yema. The web search found mainly bibliographic and catalogue pages. It found no school materials.

\section{Primary Educational Resources}
No primers, textbooks or school materials in Suena turned up in the search. The language is not known to be taught in schools.

\section{Linguistic and Reference Materials}
\begin{itemize}
\item Darryl Wilson, \emph{Suena Grammar} (1974, SIL Workpapers in Papua New Guinea Languages, vol.\ 8), the main grammatical description. It is cited in the Endangered Languages Project bibliography: \url{https://www.endangeredlanguages.com/lang/5427/bibliography}
\item A Suena--English and English--Suena dictionary compiled by Lael and Darryl Wilson is catalogued in Glottolog: \url{https://glottolog.org/resource/reference/id/311748}
\item Wilson's Suena data are cited in the Cross-Gram sources list: \url{https://crossgram.clld.org/sources/veselinovanegex-RN344}
\item Wilson's counting-system data for Suena, in the Max Planck numerals database: \url{https://lingweb.eva.mpg.de/channumerals/Suena.htm}
\item Wikipedia overview of the language: \url{https://en.wikipedia.org/wiki/Suena_language}
\item Word-list entry for the Yema-Yarawe dialect, from Smallhorn (2011), at TransNewGuinea.org: \url{https://transnewguinea.org/word/finger?page=4}
\end{itemize}

\section{Religious and Community Resources}
Global Recordings Network lists Christian materials for Suena (videos, Bibles, songs, a Jesus Film version) but reports no audio recordings currently available: \url{https://globalrecordings.net/en/language/sue}. Joshua Project also has a language page: \url{https://www.joshuaproject.net/index.php/languages/sue}.

\section{Digital and Audio Resources}
No digital corpus, audio archive or software for Suena was found beyond the catalogue entries above.

\section{Curriculum Development Context}
Little exists for building a curriculum. The Wilson grammar and dictionary are old and are held mainly in SIL channels. A curriculum would need new community-based materials, with orthography and literacy work starting from Wilson's description.
\end{document}